%% ma: L12-B08-0144
%% lop: 12
%% chuong: 2
%% bai: 8
%% dang: 12.8.49
%% loai: TN4
%% muc: VD
%% dap_an: "A"
%% nguon: "CĐ2.3.pdf giáo viên gửi 27/10/2026 (M07, Lê Minh Khôi – Nguyễn Anh Khôi, Chương II Chủ đề 3 Biểu thức tọa độ của các phép toán vectơ), Câu 60"
%% kien_thuc: "Bài 8 (điều kiện vuông góc, biểu thức tọa độ), lớp 10 (hệ phương trình bậc nhất)"
%% trang_thai: cho-duyet
%% ly_do: "Dạng mới 12.8.49 đề xuất, chờ giáo viên duyệt; Đề gốc không có phương án đúng; đã thay phương án A bằng H(1; 2/5; 4/5)"
%% ghi_chu: "Đề gốc không có phương án đúng (trực tâm đúng là (1; 2/5; 4/5)); đã thay phương án A (1; 1/7; 5/7) bằng giá trị đúng"
\begin{ex}
Trong không gian $Oxyz$, cho ba điểm $A(1;1;1)$, $B(2;3;0)$ và $C(0;2;3)$. Tọa độ trực tâm $H$ của tam giác $ABC$ là
\choice{\True $H\left(1;\dfrac25;\dfrac45\right)$}{$H\left(\dfrac17;1;\dfrac57\right)$}{$H\left(-1;-\dfrac17;-\dfrac57\right)$}{$H\left(\dfrac57;\dfrac17;1\right)$}
\loigiai{$\overrightarrow{AH}=s\overrightarrow{AB}+t\overrightarrow{AC}$ ($H$ thuộc mặt phẳng $(ABC)$), $\overrightarrow{AB}=(1;2;-1)$, $\overrightarrow{AC}=(-1;1;2)$, $\overrightarrow{BC}=(-2;-1;3)$.

$\overrightarrow{AH}\cdot\overrightarrow{BC}=0$, $\overrightarrow{BH}\cdot\overrightarrow{AC}=0\Rightarrow s=t=-\dfrac15$, $H\left(1;\dfrac25;\dfrac45\right)$.}
\end{ex}
